% SOURCE_STATEMENT: "Give a mathematically rigorous justification of the techniques used to determine molecular configurations from first principles."
% Preparation and final self-audit: CONTRIBUTING.md prompts C and D; source/audit artifacts retained in the local physics_statements_2026_10_11 batch.
% !TeX program = lualatex
\documentclass[11pt]{article}
\usepackage[margin=25mm]{geometry}
\usepackage{fontspec}
\setmainfont{Latin Modern Roman}
\usepackage{amsmath,amssymb,mathtools}
\usepackage{unicode-math}
\setmathfont{Latin Modern Math}
\usepackage{xurl,hyperref}
\hypersetup{colorlinks=true,urlcolor=blue}
\setcounter{secnumdepth}{0}
\setlength{\parindent}{0pt}
\setlength{\parskip}{5pt}
\setlength{\emergencystretch}{3em}
\newcommand{\sref}[2]{\hyperlink{source:#1:#2}{[S#2]}}
\newcommand{\cataloguescope}[1]{\paragraph{Recorded target.}#1}
\newcommand{\R}{\mathbb R}
\newcommand{\C}{\mathbb C}
\newcommand{\N}{\mathbb N}
\newcommand{\Z}{\mathbb Z}
\newcommand{\Q}{\mathbb Q}
\newcommand{\F}{\mathbb F}
\newcommand{\eps}{\varepsilon}
\begin{document}
\section{First-Principles Molecular Configurations}\label{statement}

\subsection{Definitions and mathematical statement}

\paragraph{Shared notation.}$\mathbb N=\{1,2,\ldots\}$; the standard number fields have their usual meanings. All additional conventions are stated locally.


\paragraph{Definitions and notation.}A molecular configuration specifies the nuclear positions, usually modulo rigid motions and permutation of identical nuclei. In a first-principles quantum model, electronic and nuclear degrees of freedom interact through a many-body Hamiltonian. Effective energy surfaces and their minimizing configurations underlie common prediction methods. A justification must specify the Hamiltonian, approximation, parameter regime and relationship between predicted and exact configurations.

\paragraph{Catalogue wording retained for traceability.}
Give a mathematically rigorous justification of the techniques used to determine molecular configurations from first principles.
\paragraph{Scope and source qualification.}The question does not pick one approximation scheme or completion theorem. The definitions explain the programme and do not restrict it silently to the Born--Oppenheimer approximation.
\subsection{Short English statement}
Justify mathematically the first-principles methods used to predict molecular shapes.
\subsection{Sources}
\begin{itemize}
\item[\hypertarget{source:statement:1}{S1}] ULAM.AI and the UnsolvedMath Contributors, UnsolvedMath, catalogue entry. Catalogue curation under CC BY 4.0; definitions and English exposition added locally.
\url{https://huggingface.co/datasets/ulamai/UnsolvedMath}.
\item[\hypertarget{source:statement:2}{S2}] Barry Simon, Schrödinger operators in the twenty-first century, Mathematical Physics 2000, pp. 283–288, original Problem 12.
\url{https://web.ma.utexas.edu/mp_arc-bin/mpa?yn=00-78}.
\item[\hypertarget{source:statement:3}{S3}] Wikipedia contributors, Simon problems, pinned revision 1335881350, The 2000 list. Existing catalogue wording adapted; CC BY-SA 4.0.
\url{https://en.wikipedia.org/w/index.php?title=Simon_problems&oldid=1335881350#The_2000_list}.
\end{itemize}
\label{end:statement}
\end{document}
